\subsection{A smooth category retaining the actual collision corner}
The direct smoothness failure in \fid{F39} has an explicit repair
at the level of a larger category. The following is an application
of the classical Auslander construction to the actual algebra and
module already identified in \fid{F19}/\fid{F21}; no priority claim
for the Auslander algebra is intended.

\begin{finding}[\fid{F40}: a smooth proper corner and an exact module lift]\label{f:auslander}
Let $k=\F_q$ have characteristic greater than three, and put
\[
 B=k[\epsilon]/(\epsilon^2),\qquad
 \Gamma=\End_B(B\oplus k),\qquad
 \Lambda=k\times k\times M_2(\Gamma).
\]
Then $\dim\Gamma=5$, $\dim\Lambda=22$, and both algebras are
smooth proper of global dimension two. For an explicit idempotent
$f\in\Lambda$,
\[
 f\Lambda f=A_1=k\times k\times M_2(B).
\]
There are projective right $\Lambda$-modules
$\mathcal P=f\Lambda$ and $\widetilde M$, with
$\dim_k\widetilde M=104$, such that
\[
 \End_\Lambda(\mathcal P)=A_1,\qquad
 \operatorname{Hom}_\Lambda(\mathcal P,\widetilde M)
   \simeq M_1
\]
as the actual right $A_1$-module of dimension 64 in \fid{F21}.
Thus the full algebra, including its nilpotents, and the specified
specialized tensor action are recovered.

The induced functor
$\operatorname{Perf}(A_1)\to\operatorname{Perf}(\Lambda)$ is fully
faithful but not right-admissible. The noncommutative Chow motive
of $\Lambda$ is a sum of four unit motives. Its normalized
cyclotomic Frobenius therefore has four eigenvalues one on
$TP_0[1/p]$, and $TP_1[1/p]=0$.
No relative deformation or based-operation comparison is asserted.
\end{finding}

\begin{proof}
Let $e_1,e_2$ be the idempotents of the summands $B,k$.
Let $a:k\to B$ send $1$ to $\epsilon$, let $b:B\to k$ be
the quotient, and let $c=ab$. Multiplication is composition,
with the right factor applied first. The five elements
$e_1,e_2,a,b,c$ form a basis and satisfy
\[
 a=e_1ae_2,\qquad b=e_2be_1,\qquad ba=0,\qquad
 c=ab\ne0,\qquad c^2=0.
\]
Indeed the four matrix Hom spaces have dimensions $2,1,1,1$,
and these maps span them. The radical is
$J=\langle a,b,c\rangle$, with $J^2=kc$, $J^3=0$ and
$\Gamma/J=k^2$.

Write
\[
 \mathsf P_1=e_1\Gamma=\langle e_1,a,c\rangle,\qquad
 \mathsf P_2=e_2\Gamma=\langle e_2,b\rangle
\]
and let $S_i$ be their simple tops.
Left multiplication by $a$ sends $e_2,b$ to $a,c$, and left
multiplication by $b$ sends $e_1$ to $b$ while killing $a,c$.
Thus the following are exact projective resolutions:
\[
 0\longrightarrow\mathsf P_2\xrightarrow{a}\mathsf P_1
   \longrightarrow S_1\longrightarrow0,
\]
\begin{equation}\label{eq:auslandersimple}
 0\longrightarrow\mathsf P_2\xrightarrow{a}\mathsf P_1
   \xrightarrow{b}\mathsf P_2\longrightarrow S_2
   \longrightarrow0.
\end{equation}
Every right module has a radical filtration of length at most
three, with semisimple quotients. Direct sums of these resolutions
and the extension bound for projective dimension show that every
module has projective dimension at most two. Applying
$\operatorname{Hom}_\Gamma(-,S_2)$ to
\eqref{eq:auslandersimple} gives
$\operatorname{Ext}^2_\Gamma(S_2,S_2)=k$, since
$\operatorname{Hom}(\mathsf P_1,S_2)=0$ and
$\operatorname{Hom}(\mathsf P_2,S_2)=k$. Hence the global
dimension is exactly two.

We prove smoothness directly. All tensor products in the following
diagonal resolution are over $k$. Set
\[
 Q_0=(\Gamma e_1\otimes e_1\Gamma)
          \oplus(\Gamma e_2\otimes e_2\Gamma),
\]
\[
 Q_1=(\Gamma e_1\otimes e_2\Gamma)_a
          \oplus(\Gamma e_2\otimes e_1\Gamma)_b,\qquad
 Q_2=\Gamma e_2\otimes e_2\Gamma.
\]
These are projective $\Gamma$-bimodules. Let $\mu:Q_0\to\Gamma$
be multiplication and define
\[
 d_1(x\otimes y)_a=xa\otimes y-x\otimes ay,\qquad
 d_1(x\otimes y)_b=xb\otimes y-x\otimes by,
\]
\[
 d_2(x\otimes y)=(xb\otimes y)_a+(x\otimes ay)_b.
\]
The subscripts specify the summands; the terms of $d_1$ lie in
the corresponding source and target vertex summands of $Q_0$.
These are bimodule maps, with $\mu d_1=0$ and
\[
 d_1d_2(x\otimes y)=xba\otimes y-x\otimes bay=0.
\]
The left ideals have bases
$\Gamma e_1=\langle e_1,b,c\rangle$ and
$\Gamma e_2=\langle e_2,a\rangle$. Consequently
\[
 \dim Q_2=4,\quad\dim Q_1=12,\quad\dim Q_0=13,\quad
 \dim\Gamma=5.
\]
The map $\mu$ is surjective. The image of $d_1$ contains the
eight independent vectors
\[
 \begin{gathered}
 a\otimes e_2-e_1\otimes a,\quad
 a\otimes b-e_1\otimes c,\quad
 b\otimes e_1-e_2\otimes b,\quad
 c\otimes e_1-a\otimes b,\\
 b\otimes a,\quad b\otimes c,\quad
 c\otimes a,\quad c\otimes c.
 \end{gathered}
\]
Independence follows from the tensor bases: the first four have
unique coordinates at $a\otimes e_2,e_1\otimes c,e_2\otimes b,
c\otimes e_1$, respectively, and the last four are distinct
remaining coordinates. Thus $\operatorname{im}d_1=\ker\mu$.
The $a$-component of $d_2$ is the tensor product of the injection
$x\mapsto xb:\Gamma e_2\to\Gamma e_1$, which sends $e_2,a$
to $b,c$, with the identity of $e_2\Gamma$. Hence $d_2$ is
injective. Its image has dimension four, equal to
$\dim\ker d_1$, proving exactness of
\begin{equation}\label{eq:auslanderdiagonal}
 0\longrightarrow Q_2\xrightarrow{d_2}Q_1
   \xrightarrow{d_1}Q_0\xrightarrow{\mu}\Gamma
   \longrightarrow0.
\end{equation}
This is a finite projective diagonal resolution. It proves
homological smoothness; finite dimensionality proves properness.
Matrix Morita equivalence and finite products give the same
properties and global dimension for $\Lambda$.

Now put $f=(1,1,I_2\otimes e_1)$. The corner equality follows
from $e_1\Gamma e_1=\langle e_1,c\rangle\simeq B$, with
$\epsilon\mapsto c\ne0$. In particular all named generators in
\fid{F19} are retained: its normalized matrix formulas use
$z=3/4+c$ at the collision. Left multiplication identifies
$\End_\Lambda(f\Lambda)=f\Lambda f=A_1$.

For the module recovery, write the actual specialization from
\fid{F21} as
\[
 M_1=\epsilon_+^{20}\oplus\epsilon_-^4\oplus S^{20},
\]
where $S=k^2$ is the right row module for $M_2(B)$, with
$\epsilon$ acting as zero. The projective right $M_2(\Gamma)$-module
\[
 Q=(E_{11}\otimes e_2)M_2(\Gamma)
     \simeq(\mathsf P_2,\mathsf P_2)
\]
has dimension four. Its corner is
\[
 Q(I_2\otimes e_1)=(\mathsf P_2e_1)^2
       =(kb)^2\simeq S,
\]
as an actual module, since $bc=bab=0$. Therefore
\[
 \widetilde M=\epsilon_+^{20}\oplus\epsilon_-^4\oplus Q^{20}
\]
is projective of dimension $104$, and
$\widetilde M f=M_1$. Evaluation at $f$ identifies
$\operatorname{Hom}_\Lambda(f\Lambda,\widetilde M)$ with this
corner. Precomposition by $\End_\Lambda(f\Lambda)$ is precisely
the right $A_1$-action, proving the claimed action-preserving
comparison.

For the embedding and its limitation, first consider
\[
 I:\operatorname{Perf}(B)\longrightarrow\operatorname{Perf}(\Gamma),
 \qquad X\longmapsto X\otimes_B^{\mathbf L}e_1\Gamma.
\]
It sends the generator $B$ to the projective $\mathsf P_1$.
Its right adjoint on ambient derived categories is corner
restriction $Y\mapsto Ye_1$. The unit is an isomorphism because
\[
 (X\otimes_B^{\mathbf L}e_1\Gamma)e_1
       =X\otimes_B^{\mathbf L}B=X.
\]
Thus $I$ is fully faithful. The matrix and scalar factors give
the assertion for $A_1,\Lambda$.

However $\mathsf P_2e_1=kb\simeq k$ is not perfect over $B$:
the exact resolution
\[
 \cdots\xrightarrow{\epsilon}B\xrightarrow{\epsilon}B
   \xrightarrow{\epsilon}B\longrightarrow k\longrightarrow0
\]
gives $\operatorname{Tor}^B_i(k,k)=k$ in every degree.
If $I$ had a right adjoint valued in $\operatorname{Perf}(B)$,
its value at $\mathsf P_2$ would represent
$X\mapsto\operatorname{Hom}_{D(B)}(X,k)$ on perfect $X$.
Evaluation at $B[n]$ would force the representing object to
have only the cohomology module $k$ in degree zero; naturality
for $\End_B(B)=B$ fixes its $B$-action. It would be isomorphic
to the nonperfect object $k$, a contradiction. This also
obstructs right admissibility of the actual matrix corner.

Finally $\langle S_1,\mathsf P_2\rangle$ is a full exceptional
pair in $\operatorname{Perf}(\Gamma)$.
The projective $\mathsf P_2$ has endomorphism ring $k$.
The first simple resolution gives
$\End(S_1)=k$ and no positive self-Ext, because
$\operatorname{Hom}(\mathsf P_2,S_1)=S_1e_2=0$.
Also $R\operatorname{Hom}(\mathsf P_2,S_1)=0$, so the pair
is semiorthogonal. Its triangle recovers $\mathsf P_1$, hence
both projective generators. Each exceptional subcategory has
both adjoints given by finite derived Hom and its dual, so this
is a semiorthogonal decomposition into two copies of
$\operatorname{Perf}(k)$.

Known semiorthogonal additivity
\cite[Example~12.3, p.~24]{TabuadaWeil} now gives
\[
 U(\Gamma)_{\Q}=U(k)_{\Q}^{\oplus2},\qquad
 U(\Lambda)_{\Q}=U(k)_{\Q}^{\oplus4}.
\]
The Frobenius-enriched TP functor
\cite[\S6.2, equation~(6.9), p.~16]{TabuadaWeil}
preserves these isomorphisms with Frobenius. The point comparison
\cite[\S7, equations~(7.1)--(7.4), p.~17]{TabuadaWeil}
then gives $TP_0(\Lambda)[1/p]=K^4$ with $F_0=1$ and
$TP_1(\Lambda)[1/p]=0$, where $K=\operatorname{Frac}W(k)$.
This is also an instance of the cited connective smooth proper
theorem and its motive computation
\cite[Theorem~2.12, p.~8; proof, equation~(12.6), p.~26]{TabuadaWeil}.
The source hypotheses hold by \eqref{eq:auslanderdiagonal} and
finite dimensionality. The special-case exceptional calculation
above avoids any assumption of nilinvariance for arbitrary
nonsmooth algebras.
\end{proof}

The original $M_1$ is not perfect over $A_1$. Its recovery here
is a corner restriction of a projective object, not the image of
$M_1$ under the perfect-module embedding. Its higher coherent
self-Ext need not agree with that of $\widetilde M$, and the
relative collision order in \fid{F25} has not been recovered
from this single fiber. The pure spectrum of the entire smooth
category does not identify a Frobenius realization of every
specified corner morphism, a positive nonstandard basis, or its
structure constants. The direct smoothness gate is nevertheless
strictly narrowed: this category retains the actual collision
algebra and tensor action as stated, without semisimplifying them.

% Separate paragraph for Part III.
\paragraph{\fid{F40} in Grothendieck's language.}
Affine endomorphism algebras, idempotent corners and perfect
diagonals give the precise coherent dictionary. The explicit
finite diagonal resolution proves smoothness of a larger category,
while its projective objects retain the singular corner and its
actual module. This uses the perfect-complex language associated
with SGA~6 and the explicitly cited modern smooth dg and
noncommutative-motive theory. The exceptional decomposition
explains the unit Frobenius spectrum without turning the corner
subcategory into an admissible motivic summand. No unpublished
archive leaf is claimed to contain this construction. A relative
order, a realization of the designated morphisms and a comparison
of the required bases remain separate mathematical tasks.
